% !TeX program = lualatex
% ============================================================
% chapter3.tex — ch03 唯一內容源（2026-08-09 起；LaTeX 統一 U1）
% 沿革：升格自 dist/ch03/chapter3.tex（原 make_dist.py 自 fragment 轉換的成品；
%   fragment 樹已凍結，改課文一律改本檔。拍板見 ../../KICKOFF-latex-unification.md）
% 樣式源＝../../template/calcbook.sty（M-B1 拍板模板）
% 編譯（本資料夾為 CWD；aux 進 build/、成品 PDF 移入 ../../dist/ch03/）：
%   latexmk -lualatex -auxdir=../../build/aux-ch03 chapter3.tex
% ============================================================
\documentclass[a4paper,12pt,oneside]{memoir}
\makeatletter\def\input@path{{../../template/}}\makeatother
\def\cbfontsdir{../../template/fonts/inter/}
\usepackage{calcbook}
% 圖：emitter 射 `<ch>/<stem>`（見 convert.py），故 graphicspath 指向 chapters/ 上層。
% 模板內的 {../figs/} 是 chapters/<ch>/ 為 CWD 時的路徑，dist/<ch>/ 需這一行覆寫。
\graphicspath{{../../chapters/}}
\begin{document}
\cbchapter{3}
\chapteropener{Chapter 3}{Chain Rule and Trigonometric Derivatives}
\begin{lead}
Chapter 2 built the derivative from its definition as the limit of a difference quotient, and turned it into a working toolkit: rules for constants, powers, sums, products, quotients, and the exponential. Yet a whole world of functions still lies just out of reach. How fast does \(\sin x\) change? And what is the derivative of a \emph{composition} such as \(\sqrt{1 + x^{2}}\) or \(\sin(x^{2})\), where the output of one function becomes the input of another? This chapter supplies the two ideas that complete the differentiation toolkit. First we differentiate the trigonometric functions. Algebra alone cannot evaluate the limit that this requires, so we use a geometric comparison of areas on the unit circle. Then we prove the \emph{chain rule}, the single rule that differentiates any composition, and use it to unlock the derivatives of inverse and implicitly defined functions.
\end{lead}
\parahead{By the end of this chapter, you will be able to:}
\begin{objectives}
  \item differentiate \(\sin x\), \(\cos x\), \(\tan x\), and the other trigonometric functions, and combine them with the rules of Chapter 2;
  \item state the chain rule and use it to differentiate compositions \(f(g(x))\), including nested ones;
  \item apply logarithmic differentiation to handle expressions such as \(x^{x}\) and \(f(x)^{g(x)}\);
  \item derive the derivatives of \(\arcsin x\), \(\arctan x\), and \(\ln x\) from their inverse-function relations.
\end{objectives}
\sechead{3.1}{Derivatives of the Sine and Cosine Functions}
So far algebra has been enough to compute every derivative. For a polynomial we expanded \((x + h)^{n}\) and canceled the factor of \(h\); for \(e^{x}\) the power series did the same job. The trigonometric functions are different in kind. Consider the difference quotient for \(\sin x\), \[ \frac{\sin(x + h) - \sin x}{h}. \] No algebraic identity cancels the \(h\) in the denominator against the numerator. Expanding \(\sin(x + h) - \sin x\) does not produce a factor of \(h\). As we are about to see, the entire computation reduces to a single limit, \[ \lim_{\theta \to 0} \frac{\sin \theta}{\theta}, \] which has the indeterminate form \(0/0\) and which no amount of algebra will resolve. To evaluate it we need a genuinely new tool: a geometric comparison of areas on the unit circle, squeezed to a limit. That one limit, once secured, is the main step toward the derivatives of \emph{both} \(\sin\) and \(\cos\), and through them the rest of the trigonometric functions.

Throughout this section angles are measured in radians; this assumption turns out to be essential, as we note below.

\subsechead{The difference quotient for sine}
Fix \(x\) and form the difference quotient for \(f(x) = \sin x\). The decisive algebraic step is the sum-to-product identity \(\sin A - \sin B = 2 \cos\frac{A + B}{2} \sin\frac{A - B}{2}\). If this identity is not familiar, it follows from the angle-sum formulas. Set \(u = \frac{A + B}{2}\) and \(v = \frac{A - B}{2}\), so that \(A = u + v\) and \(B = u - v\). Expanding \(\sin(u + v) - \sin(u - v)\) by the angle-sum formulas cancels the \(\sin u \cos v\) terms and leaves exactly \(2 \cos u \sin v\). (Appendix A.2 derives the whole family of such identities.) Applied with \(A = x + h\) and \(B = x\): \[ \sin(x + h) - \sin x = 2 \cos\Bigl(x + \tfrac{h}{2}\Bigr) \sin\frac{h}{2}. \] Dividing by \(h\) and writing the denominator as \(h = 2 \cdot \tfrac{h}{2}\), to match the \(\tfrac{h}{2}\) inside the sine, gives

\[ \frac{\sin(x + h) - \sin x}{h} = \cos\Bigl(x + \tfrac{h}{2}\Bigr) \cdot \frac{\sin(h/2)}{h/2}. \]

As \(h \to 0\), two things happen at once. The factor \(\cos\bigl(x + \tfrac{h}{2}\bigr)\) ought to tend to \(\cos x\) — provided \(\cos\) is continuous — and the factor \(\dfrac{\sin(h/2)}{h/2}\) is exactly \(\dfrac{\sin\theta}{\theta}\) with \(\theta = h/2 \to 0\). So the derivative rests on two facts: that sine and cosine are continuous, and the value of \(\lim_{\theta \to 0}\sin\theta/\theta\). We secure both now.

\subsechead{A geometric inequality on the unit circle}
To get a grip on \(\sin\theta/\theta\) for small \(\theta\), we compare three areas on a circle of radius \(1\). Because both \(\theta\) and \(\sin\theta\) reverse sign when \(\theta\) is replaced by \(-\theta\), the ratio \(\sin\theta/\theta\) is unchanged. A function with this property is called an \emph{even} function. It is therefore enough to treat \(\theta > 0\), shrinking to \(0\). Take \(0 < \theta < \tfrac{\pi}{2}\) and look at Figure \ref{fig:3.1}.

\begin{figureblock}
  \includegraphics[width=79.9mm]{ch03/figs/sector-inequality}
\figcaption{fig:3.1}{The unit-circle area comparison behind \(\sin\theta \le \theta \le \tan\theta\).}
\end{figureblock}
Here \(A = (1, 0)\) lies on the circle, \(B = (\cos\theta, \sin\theta)\) is the point at angle \(\theta\), and \(C = (1, \tan\theta)\) is where the radius \(OB\), extended, meets the vertical tangent line at \(A\). Three nested regions appear, and their areas are easy to read off:

\begin{itemize}
  \item the inscribed triangle \(OAB\) has base \(OA = 1\) and height \(\sin\theta\), so its area is \(\tfrac12 \sin\theta\);
  \item the circular sector \(OAB\) is the fraction \(\theta / (2\pi)\) of a disc of area \(\pi\), so its area is \(\tfrac12 \theta\);
  \item the larger triangle \(OAC\) has base \(OA = 1\) and height \(AC = \tan\theta\), so its area is \(\tfrac12 \tan\theta\).
\end{itemize}
The triangle \(OAB\) lies inside the sector, which lies inside the triangle \(OAC\), so their areas increase in that order: \[ \tfrac12 \sin\theta \;\le\; \tfrac12 \theta \;\le\; \tfrac12 \tan\theta. \] Equivalently, split the outer triangle as \(OAC = OAB \cup ABC\): the corner piece \(ABC\) is what \(OAC\) adds to the inscribed triangle \(OAB\), and it contains the circular segment between the chord \(AB\) and the arc. Multiplying through by \(2\) gives \(\sin\theta \le \theta \le \tan\theta\). Dividing by \(\sin\theta > 0\) and using \(\tan\theta = \sin\theta/\cos\theta\) turns this into \(1 \le \dfrac{\theta}{\sin\theta} \le \dfrac{1}{\cos\theta}\), and taking reciprocals of these positive quantities reverses the inequalities: \[ \cos\theta \;\le\; \frac{\sin\theta}{\theta} \;\le\; 1, \qquad 0 < \theta < \tfrac{\pi}{2}. \tag{1} \] In particular \(0 \le \lvert \sin\theta \rvert \le \lvert \theta \rvert\) for \(\lvert \theta \rvert < \tfrac{\pi}{2}\), a bound we use repeatedly below.

\subsechead{Continuity of sine and cosine}
The bound \(\lvert \sin\theta \rvert \le \lvert \theta \rvert\) does more than help with the fundamental limit: it is also what makes sine and cosine continuous, in the sense introduced in Chapter 1. We need continuity for the first factor of the difference quotient and for the proof of the fundamental limit, so we establish it now.

\begin{envproposition}{Proposition}{prop:3.1}{Continuity of sine and cosine}
The functions \(\sin x\) and \(\cos x\) are continuous at every \(x_{0} \in \mathbb{R}\).
\end{envproposition}
\begin{envproof}{Proof}{}{}
Everything rests on one bound: \(0 \le \lvert \sin\theta \rvert \le \lvert \theta \rvert\) near \(\theta = 0\), from inequality (1). Since \(\lvert \theta \rvert \to 0\) as \(\theta \to 0\), the squeeze theorem (§1.5) forces \(\lvert \sin\theta \rvert \to 0\), so

\[ \lim_{\theta \to 0} \sin\theta = 0. \]

Now fix \(x_{0}\). The sum-to-product identities give

\[ \begin{aligned}
        \cos x - \cos x_{0} &= -2 \sin\frac{x - x_{0}}{2}\, \sin\frac{x + x_{0}}{2}, \\
        \sin x - \sin x_{0} &= \phantom{-}2 \cos\frac{x + x_{0}}{2}\, \sin\frac{x - x_{0}}{2}.
      \end{aligned} \]

In each line the factor with argument \(\tfrac{x + x_{0}}{2}\) has absolute value at most \(1\), so

\[ \lvert \cos x - \cos x_{0} \rvert \le 2\Bigl\lvert \sin\tfrac{x - x_{0}}{2} \Bigr\rvert, \qquad \lvert \sin x - \sin x_{0} \rvert \le 2\Bigl\lvert \sin\tfrac{x - x_{0}}{2} \Bigr\rvert. \]

As \(x \to x_{0}\) the angle \(\tfrac{x - x_{0}}{2} \to 0\), so by the limit just established both right-hand sides tend to \(0\). The squeeze theorem then gives \(\cos x \to \cos x_{0}\) and \(\sin x \to \sin x_{0}\); that is, cosine and sine are continuous at \(x_{0}\).
\end{envproof}
\subsechead{The fundamental limit}
With continuity secured, we can finally squeeze \(\sin\theta/\theta\) to its limit (Figure \ref{fig:3.2}).

\begin{envproposition}{Proposition}{prop:3.2}{The fundamental trigonometric limit}
\[ \lim_{\theta \to 0} \frac{\sin\theta}{\theta} = 1. \]
\end{envproposition}
\begin{envproof}{Proof}{}{}
For \(0 < \theta < \tfrac{\pi}{2}\), inequality (1) reads \(\cos\theta \le \dfrac{\sin\theta}{\theta} \le 1\). As \(\theta \to 0^{+}\) the upper bound is the constant \(1\), while the lower bound satisfies \(\cos\theta \to \cos 0 = 1\) by the continuity of cosine (Proposition \ref{prop:3.1}). The squeeze theorem (§1.5) pins the quantity between them: \(\lim_{\theta \to 0^{+}} \sin\theta/\theta = 1\). Finally \(\sin\theta/\theta\) is even, so its left-hand limit at \(0\) equals its right-hand limit, and the two-sided limit is \(1\) as well.
\end{envproof}
\begin{figureblock}
  \includegraphics[width=81.92mm]{ch03/figs/squeeze-limit}
\figcaption{fig:3.2}{A graphical reading of the squeeze in Proposition \ref{prop:3.2}: \(\dfrac{\sin\theta}{\theta}\) sits between \(\cos\theta\) and \(1\), with an open circle at \(\theta = 0\).}
\end{figureblock}
\begin{envcaution}{Caution}{}{}
Proposition \ref{prop:3.2} is a statement about the \emph{limit} as \(\theta \to 0\), not an identity. At a fixed angle the ratio \(\sin\theta/\theta\) is generally not \(1\): at \(\theta = \tfrac{\pi}{2}\) it equals \(\dfrac{\sin(\pi/2)}{\pi/2} = \dfrac{1}{\pi/2} = \dfrac{2}{\pi} \approx 0.64\), and at \(\theta = \pi\) it equals \(\dfrac{0}{\pi} = 0\). The value \(1\) is reached only in the limit, as \(\theta\) shrinks to \(0\).
\end{envcaution}
\begin{envcaution}{Caution}{}{}
Everything here depends on measuring angles in \emph{radians}. The formula \(\tfrac12\theta\) for the sector area uses radian measure, for which the arc cut off on the unit circle has length exactly \(\theta\). The limit \(\sin\theta/\theta \to 1\) rests on that formula, so it depends on radians too. In degrees the limit is not \(1\): since \(x^{\circ} = \tfrac{\pi}{180}x\) radians, one finds \(\lim_{x \to 0}\dfrac{\sin(x^{\circ})}{x} = \dfrac{\pi}{180}\), and so \(\dfrac{d}{dx}\sin(x^{\circ}) = \dfrac{\pi}{180}\cos(x^{\circ})\). The clean formulas of this section, \(\sin' = \cos\) and \(\cos' = -\sin\), hold only in radians.
\end{envcaution}
\subsechead{The derivatives of sine and cosine}
With continuity and the fundamental limit established, the difference quotient for \(\sin\) is now easy to evaluate.

\begin{envtheorem}{Theorem}{thm:3.1}{Derivative of sine}
\[ \frac{d}{dx}\sin x = \cos x \qquad \text{for all } x \in \mathbb{R}. \]
\end{envtheorem}
\begin{envproof}{Proof}{}{}
Using the form of the difference quotient found above,

\[ \begin{aligned}
        \frac{d}{dx}\sin x &= \lim_{h \to 0} \frac{\sin(x + h) - \sin x}{h} \\[2pt]
          &= \lim_{h \to 0} \cos\Bigl(x + \tfrac{h}{2}\Bigr) \cdot \frac{\sin(h/2)}{h/2}.
      \end{aligned} \]

As \(h \to 0\), the first factor tends to \(\cos x\) because cosine is continuous (Proposition \ref{prop:3.1}), and the second tends to \(1\) by the fundamental limit (Proposition \ref{prop:3.2}), with \(\theta = h/2 \to 0\). The product therefore tends to \(\cos x \cdot 1 = \cos x\).
\end{envproof}
The same machinery handles cosine, this time through the companion identity \(\cos A - \cos B = -2 \sin\frac{A + B}{2} \sin\frac{A - B}{2}\), obtained just like the sine version by expanding \(\cos(u + v) - \cos(u - v)\) with the angle-sum formulas.

\begin{envtheorem}{Theorem}{thm:3.2}{Derivative of cosine}
\[ \frac{d}{dx}\cos x = -\sin x \qquad \text{for all } x \in \mathbb{R}. \]
\end{envtheorem}
\begin{envproof}{Proof}{}{}
With \(A = x + h\) and \(B = x\), the identity gives \(\cos(x + h) - \cos x = -2 \sin\bigl(x + \tfrac{h}{2}\bigr) \sin\tfrac{h}{2}\), so

\[ \frac{\cos(x + h) - \cos x}{h} = -\sin\Bigl(x + \tfrac{h}{2}\Bigr) \cdot \frac{\sin(h/2)}{h/2}. \]

Letting \(h \to 0\), the factor \(\sin\bigl(x + \tfrac{h}{2}\bigr) \to \sin x\) by continuity (Proposition \ref{prop:3.1}) and \(\dfrac{\sin(h/2)}{h/2} \to 1\) by Proposition \ref{prop:3.2}, so

\[ \frac{d}{dx}\cos x = -\sin x \cdot 1 = -\sin x. \]
\end{envproof}
\begin{figureblock}
  \includegraphics[width=55.07mm]{ch03/figs/sin-cos-slope}
\figcaption{fig:3.3}{Tangent slopes on \(y = \sin x\) matching the heights of \(y = \cos x\) at the same \(x\)-values.}
\end{figureblock}
\begin{workedexample}
\begin{envexample}{Example}{ex:3.1}{}
Establish the companion limit \(\displaystyle \lim_{\theta \to 0} \frac{1 - \cos\theta}{\theta} = 0\).
\end{envexample}
\begin{envsolution}{Solution}{}{}
This is again of the form \(0/0\). Multiplying numerator and denominator by \(1 + \cos\theta\) turns the numerator into \(1 - \cos^{2}\theta = \sin^{2}\theta\):

\[ \begin{aligned}
          \frac{1 - \cos\theta}{\theta}
            &= \frac{(1 - \cos\theta)(1 + \cos\theta)}{\theta\,(1 + \cos\theta)}
             = \frac{1 - \cos^{2}\theta}{\theta\,(1 + \cos\theta)} \\[2pt]
            &= \frac{\sin^{2}\theta}{\theta\,(1 + \cos\theta)}
             = \frac{\sin\theta}{\theta} \cdot \frac{\sin\theta}{1 + \cos\theta}.
        \end{aligned} \]

As \(\theta \to 0\), the first factor tends to \(1\) by Proposition \ref{prop:3.2}, while the second tends to \(\dfrac{0}{1 + 1} = 0\) by the continuity of sine and cosine (Proposition \ref{prop:3.1}). The product tends to \(1 \cdot 0 = 0\).
\end{envsolution}
\end{workedexample}
\begin{workedexample}
\begin{envexample}{Example}{ex:3.2}{}
Using \(\sin' = \cos\) and \(\cos' = -\sin\) together with the quotient rule (§2.5), differentiate (a) \(\tan x\) and (b) \(\sec x\), wherever \(\cos x \ne 0\).
\end{envexample}
\begin{envsolution}{Solution}{}{}
For (a), write \(\tan x = \dfrac{\sin x}{\cos x}\) and apply the quotient rule:

\[ \begin{aligned}
          \frac{d}{dx}\tan x
            &= \frac{(\cos x)(\cos x) - (\sin x)(-\sin x)}{\cos^{2} x} \\[2pt]
            &= \frac{\cos^{2} x + \sin^{2} x}{\cos^{2} x}
             = \frac{1}{\cos^{2} x} = \sec^{2} x,
        \end{aligned} \]

where the Pythagorean identity \(\cos^{2} x + \sin^{2} x = 1\) reduces the numerator to \(1\).

For (b), write \(\sec x = \dfrac{1}{\cos x}\). The quotient rule, with the constant numerator \(1\), gives

\[ \begin{aligned}
          \frac{d}{dx}\sec x
            &= \frac{0 \cdot \cos x - 1 \cdot (-\sin x)}{\cos^{2} x} \\[2pt]
            &= \frac{\sin x}{\cos^{2} x}
             = \frac{1}{\cos x}\cdot\frac{\sin x}{\cos x}
             = \sec x \tan x.
        \end{aligned} \]

The same technique applied to \(\cot x = \dfrac{\cos x}{\sin x}\) and \(\csc x = \dfrac{1}{\sin x}\) gives \(\dfrac{d}{dx}\cot x = -\csc^{2} x\) and \(\dfrac{d}{dx}\csc x = -\csc x \cot x\) (where \(\sin x \ne 0\)), completing the derivatives of all six trigonometric functions.
\end{envsolution}
\end{workedexample}
\begin{workedexample}
\begin{envexample}{Example}{ex:3.3}{}
A weight bobbing on a spring has height \(s(t) = \sin t\) above its rest level at time \(t\). Find its velocity \(s'(t)\) and acceleration \(s''(t)\), and describe how the acceleration is related to the height.
\end{envexample}
\begin{envsolution}{Solution}{}{}
Velocity is the derivative of height, and acceleration the derivative of velocity. Differentiating with Theorems \ref{thm:3.1} and \ref{thm:3.2},

\[ s'(t) = \cos t, \qquad s''(t) = \frac{d}{dt}\cos t = -\sin t. \]

Thus \(s''(t) = -\sin t = -s(t)\): at every instant the weight accelerates back toward its rest level, by an amount proportional to its distance from that level. This relation \(s'' = -s\) is the signature of \emph{simple harmonic motion} (Figure \ref{fig:3.4}). Both sine and cosine satisfy this relation, so both can be used to describe oscillation.
\end{envsolution}
\end{workedexample}
\begin{figureblock}
  \includegraphics[width=79.9mm]{ch03/figs/shm-triple}
\figcaption{fig:3.4}{Simple harmonic motion over one time axis: \(s = \sin t\), \(s' = \cos t\), and \(s'' = -\sin t\).}
\end{figureblock}
\begin{envremark}{Remark}{rem:3.1}{The derivative cycle of sine and cosine}
Differentiation sends sine and cosine into each other, with a sign that changes every second step (Figure \ref{fig:3.3}). Writing \(\to\) for one application of \(\dfrac{d}{dx}\),

\[ \sin x \;\to\; \cos x \;\to\; -\sin x \;\to\; -\cos x \;\to\; \sin x. \]

After four derivatives we are back where we started, so \(\dfrac{d^{4}}{dx^{4}}\sin x = \sin x\), and likewise for cosine. Where \(e^{x}\) reproduces itself in a single step (§2.4), sine and cosine need four. Two of those steps already turn each of them into its own negative, which is the relation \(s'' = -s\) of Example \ref{ex:3.3}.
\end{envremark}
We can now differentiate every trigonometric function, but only in its bare form, \(\sin x\) rather than \(\sin(x^{2})\) or \(\sin(3x + 1)\). To reach a function tucked \emph{inside} another we need a rule for differentiating compositions. That rule, the chain rule, is the business of the next section; with it, the derivatives found here become the seeds of a vast family of new ones.


\sechead{3.2}{The Chain Rule}
Section 3.1 gave us the derivatives of the trigonometric functions, but only in their bare form. We can differentiate \(\sin x\), yet nothing so far reaches \(\sin(x^{2})\) or \(\sqrt{1 + x^{2}}\), where one function is fed \emph{into} another. Chapter 2 supplied rules for sums, products, and quotients, but not for \emph{composition}, the operation of applying one function to the output of another. The chain rule fills exactly this gap, and it is the rule that finally completes the differentiation toolkit.

The idea behind it, pictured in Figure \ref{fig:3.5}, fits in a single image. Near a point, a differentiable function behaves like its tangent line: it takes a small change in its input and scales it by its slope. The catch is only that this is approximate: there is an error, and it shrinks faster than the change itself. Now chain two such functions together. If \(g\) scales a small increment \(h\) by \(g'(x)\), and \(f\) then scales \emph{that} by \(f'(g(x))\), the two magnifications stack. The composite scales \(h\) by the product \(f'(g(x))\,g'(x)\). That product is the chain rule, and the proof at the end of this section is nothing more than the careful check that the two error terms, once chained, are still negligible.

Alongside the chain rule, this section draws on two results that already belong to Chapter 2, the product rule (§2.5) and the fact that a differentiable function is continuous (§2.3, Theorem 2.1). We take both as established and use them freely, without re-proving them here.

\begin{envtheorem}{Theorem}{thm:3.3}{The chain rule}
Let \(g\) be differentiable at \(x_{0}\), and let \(f\) be differentiable at \(g(x_{0})\). Then the composition \(P = f \circ g\), given by \(P(x) = f(g(x))\), is differentiable at \(x_{0}\), and

\[ P'(x_{0}) = f'(g(x_{0}))\, g'(x_{0}). \]
\end{envtheorem}
\begin{figureblock}
  \includegraphics[width=79.9mm]{ch03/figs/composed-mapping}
\figcaption{fig:3.5}{The chain rule as a composed mapping: local slope factors accumulate along \(x \mapsto u \mapsto y\).}
\end{figureblock}
\begin{envremark}{Remark}{rem:3.2}{The Leibniz form}
Writing \(u = g(x)\) for the inner variable and \(y = f(u)\) for the output, the chain rule takes the memorable Leibniz shape

\[ \frac{dy}{dx} = \frac{dy}{du}\cdot\frac{du}{dx}. \]

This is the form most often used in practice, and it reads as “the rate of \(y\) per \(x\) equals the rate of \(y\) per \(u\) times the rate of \(u\) per \(x\).” The apparent cancellation of \(du\) is a useful memory aid, but it is \emph{not} a proof: \(\dfrac{dy}{du}\) and \(\dfrac{du}{dx}\) are limits, not fractions with a common factor to cancel. What actually makes the rule true is the argument given below, where the controlled error terms — not a literal \(du/du = 1\) — do the work.
\end{envremark}
\begin{envstrategy}{Strategy}{strat:3.1}{Differentiating a composition}
To differentiate a composite expression, work from the outside in.

\begin{steps}
  \item Identify the \emph{outermost} operation — the last thing you would carry out if you evaluated the expression at a number. Call it the outer function \(f\).
  \item Whatever that operation is applied to is the inner function \(u = g(x)\).
  \item Differentiate the outer function at the inner one, forming \(f'(g(x))\) — treat the inside as a single block and do not differentiate it yet.
  \item Multiply by the derivative of the inside: the answer is \(f'(g(x))\cdot g'(x)\).
  \item If the inside is itself a composition, repeat steps 1–4 on it. Every layer contributes one slope factor, and the factors multiply.
\end{steps}
For \(\sqrt{1 + x^{2}}\), the last operation is the square root, so \(f(u) = \sqrt{u}\) and \(g(x) = 1 + x^{2}\); for \(\sin(x^{2})\), the outer function is \(\sin\) and the inner is \(x^{2}\).
\end{envstrategy}
\subsechead{Why the rule is true}
The rule is already usable. The rest of this section proves it. On a first reading you may jump straight to the subsection “Using the rule” below and return to this proof later; none of the examples depend on it. To prove that the two slopes really do multiply, we need to handle both functions’ linear approximations at once, and for that it helps to repackage what “differentiable” means. Chapter 2 defined \(f\) to be differentiable at \(x_{0}\) when the difference quotient \(\dfrac{f(x_{0}+h) - f(x_{0})}{h}\) has a limit as \(h \to 0\) (§2.2); that limit is \(f'(x_{0})\). The following form says the same thing, but as a statement about approximation by a straight line.

\begin{envdefinition}{Definition}{def:3.1}{Differentiability, remainder form}
A function \(f\) is \emph{differentiable} at \(x_{0}\) if there is a number \(m\) and a function \(R\) such that

\[ f(x_{0} + h) = f(x_{0}) + m h + R(h), \qquad \lim_{h \to 0}\frac{R(h)}{h} = 0. \]

The number \(m\) is the derivative \(f'(x_{0})\).

\begin{informal}
Informally, \(f(x_{0}+h)\) equals its tangent-line value \(f(x_{0}) + m h\) plus an error \(R(h)\) that shrinks faster than \(h\) does. The tangent line is a good fit to first order (Figure \ref{fig:3.6}).
\end{informal}
\end{envdefinition}
\begin{figureblock}
  \includegraphics[width=65.62mm]{ch03/figs/remainder-tangent-1}\hspace{6mm}\includegraphics[width=65.62mm]{ch03/figs/remainder-tangent-2}
\figcaption{fig:3.6}{Remainder-form geometry: halving \(h\) makes the vertical error \(R(h)\) much smaller, illustrating \(R(h)/h \to 0\).}
\end{figureblock}
This is not a new notion of derivative, only a new way of writing the old one.

\begin{envproposition}{Proposition}{prop:3.3}{Equivalence of the two forms}
The limit form of differentiability (§2.2) and the remainder form (Definition \ref{def:3.1}) are equivalent: \(f\) satisfies one at \(x_{0}\) if and only if it satisfies the other, and the number \(m\) is the same in both, namely \(f'(x_{0})\).
\end{envproposition}
\begin{envproof}{Proof}{}{}
Suppose first that the limit form holds, with \(\lim_{h \to 0}\dfrac{f(x_{0}+h) - f(x_{0})}{h} = m\). Define \(R(h) = f(x_{0}+h) - f(x_{0}) - m h\). Then

\[ \frac{R(h)}{h} = \frac{f(x_{0}+h) - f(x_{0})}{h} - m \;\longrightarrow\; m - m = 0, \]

so the remainder form holds with this same \(m\). Conversely, suppose the remainder form holds: \(f(x_{0}+h) = f(x_{0}) + m h + R(h)\) with \(R(h)/h \to 0\). Dividing by \(h\),

\[ \frac{f(x_{0}+h) - f(x_{0})}{h} = m + \frac{R(h)}{h} \;\longrightarrow\; m + 0 = m, \]

so the limit form holds, again with the same \(m\). The two forms therefore describe the same property and the same derivative.
\end{envproof}
With the remainder form available, the chain rule comes out by substitution. We expand \(g\) near \(x_{0}\) and \(f\) near \(g(x_{0})\), each as a linear part plus a small remainder; substituting the expansion for \(g\) into the one for \(f\), the linear parts multiply to give \(g'(x_{0})f'(g(x_{0}))\,h\), and everything left over collects into a single remainder. The whole proof is then the one estimate that this leftover still vanishes faster than \(h\).

\begin{envproof}{Proof}{}{of Theorem \ref{thm:3.3}}
We use the remainder form of differentiability (Definition \ref{def:3.1}), which by Proposition \ref{prop:3.3} is available for any differentiable function. Because \(g\) is differentiable at \(x_{0}\),

\[ g(x_{0} + h) = g(x_{0}) + m_{1} h + R_{1}(h), \qquad m_{1} = g'(x_{0}), \quad \lim_{h \to 0}\frac{R_{1}(h)}{h} = 0. \]

Because \(f\) is differentiable at \(g(x_{0})\), we may apply the remainder form to \(f\) with the increment \(m_{1} h + R_{1}(h)\) playing the role of \(h\):

\[ \begin{aligned}
        P(x_{0}+h)
          &= f\bigl(g(x_{0}+h)\bigr) \\[2pt]
          &= f(g(x_{0})) + m_{2}\bigl[m_{1} h + R_{1}(h)\bigr] + R_{2}\bigl(m_{1} h + R_{1}(h)\bigr),
      \end{aligned} \]

where \(m_{2} = f'(g(x_{0}))\) and \(\lim_{y \to 0} R_{2}(y)/y = 0\). Collecting the term linear in \(h\) and bundling the rest into a single remainder \(R_{3}\),

\[ \begin{aligned}
        &P(x_{0}+h) = P(x_{0}) + \bigl(g'(x_{0})\,f'(g(x_{0}))\bigr) h + R_{3}(h), \\[2pt]
        &\qquad R_{3}(h) = m_{2} R_{1}(h) + R_{2}\bigl(m_{1} h + R_{1}(h)\bigr),
      \end{aligned} \]

using \(f(g(x_{0})) = P(x_{0})\) and \(m_{1} m_{2} = g'(x_{0})f'(g(x_{0})) = f'(g(x_{0}))g'(x_{0})\), the same product in either order. By the remainder form, this says precisely that \(P\) is differentiable at \(x_{0}\) with \(P'(x_{0}) = f'(g(x_{0}))g'(x_{0})\). It remains only to show that

\[ \lim_{h \to 0}\frac{R_{3}(h)}{h} = 0. \]

Split the quotient into its two pieces, \(\dfrac{R_{3}(h)}{h} = m_{2}\dfrac{R_{1}(h)}{h} + \dfrac{R_{2}(m_{1} h + R_{1}(h))}{h}\). The first piece is easy: \(m_{2}\) is a constant and \(R_{1}(h)/h \to 0\), so \(m_{2} R_{1}(h)/h \to 0\). We also note for later that \(m_{1} h + R_{1}(h) \to 0\) as \(h \to 0\), since \(m_{1} h \to 0\) and \(R_{1}(h) \to 0\) (the latter because \(R_{1}(h)/h \to 0\)).

The second piece, \(R_{2}(m_{1} h + R_{1}(h))/h\), is the delicate one. Fix \(\varepsilon > 0\).

\begin{itemize}
  \item Since \(R_{2}(y)/y \to 0\), there is a \(\delta > 0\) with \(\left\lvert R_{2}(y)/y \right\rvert < \varepsilon\) whenever \(0 < \lvert y \rvert < \delta\).
  \item Since \(m_{1} h + R_{1}(h) \to 0\), there is an \(\alpha > 0\) with \(\lvert m_{1} h + R_{1}(h) \rvert < \delta\) whenever \(0 < \lvert h \rvert < \alpha\).
\end{itemize}
Now fix \(h\) with \(0 < \lvert h \rvert < \alpha\). Suppose first that \(m_{1} h + R_{1}(h) = 0\). Putting \(y = 0\) in the remainder relation \(f(g(x_{0}) + y) = f(g(x_{0})) + m_{2} y + R_{2}(y)\) forces \(R_{2}(0) = 0\), so \(R_{2}(m_{1} h + R_{1}(h)) = R_{2}(0) = 0\) and the quotient is \(0\). Otherwise \(0 < \lvert m_{1} h + R_{1}(h) \rvert < \delta\), and we factor

\[ \frac{\lvert R_{2}(m_{1} h + R_{1}(h)) \rvert}{\lvert h \rvert} = \frac{\lvert m_{1} h + R_{1}(h) \rvert}{\lvert h \rvert} \cdot \frac{\lvert R_{2}(m_{1} h + R_{1}(h)) \rvert}{\lvert m_{1} h + R_{1}(h) \rvert}. \]

By the first bullet the second factor is less than \(\varepsilon\); and by the triangle inequality the first factor satisfies \(\dfrac{\lvert m_{1} h + R_{1}(h) \rvert}{\lvert h \rvert} = \left\lvert m_{1} + \dfrac{R_{1}(h)}{h} \right\rvert \le \lvert m_{1} \rvert + \dfrac{\lvert R_{1}(h) \rvert}{\lvert h \rvert}\). Hence, in either case,

\[ \frac{\lvert R_{2}(m_{1} h + R_{1}(h)) \rvert}{\lvert h \rvert} \le \left( \lvert m_{1} \rvert + \frac{\lvert R_{1}(h) \rvert}{\lvert h \rvert} \right)\varepsilon. \]

Finally, since \(R_{1}(h)/h \to 0\), there is an \(\alpha_{1} \in (0, \alpha)\) with \(\lvert R_{1}(h) \rvert/\lvert h \rvert < 1\) for \(0 < \lvert h \rvert < \alpha_{1}\). For such \(h\),

\[ \frac{\lvert R_{2}(m_{1} h + R_{1}(h)) \rvert}{\lvert h \rvert} < \bigl( \lvert m_{1} \rvert + 1 \bigr)\varepsilon. \]

As \(\varepsilon > 0\) was arbitrary and \(\lvert m_{1} \rvert + 1\) is a fixed constant, this second piece tends to \(0\) as well. Combining the two pieces, \(R_{3}(h)/h \to 0\), and the chain rule is proved.
\end{envproof}
\subsechead{Using the rule}
\begin{workedexample}
\begin{envexample}{Example}{ex:3.4}{}
Differentiate (a) \(\sqrt{1 + x^{2}}\) and (b) \(\sin(x^{2})\).
\end{envexample}
\begin{envsolution}{Solution}{}{}
For (a), the outermost operation is the square root, so \(f(u) = \sqrt{u} = u^{1/2}\) and \(g(x) = 1 + x^{2}\). By the power rule (Chapter 2), \(f'(u) = \tfrac12 u^{-1/2}\), and \(g'(x) = 2x\). By the chain rule,

\[ \frac{d}{dx}\sqrt{1 + x^{2}} = \frac{1}{2\sqrt{1 + x^{2}}}\cdot 2x = \frac{x}{\sqrt{1 + x^{2}}}. \]

For (b), the outer function is sine and the inner is \(x^{2}\), so with \(f'(u) = \cos u\) (Theorem \ref{thm:3.1}) and \(g'(x) = 2x\),

\[ \frac{d}{dx}\sin(x^{2}) = \cos(x^{2})\cdot 2x = 2x\cos(x^{2}). \]

The factor \(2x\) is the inner derivative; dropping it and writing just \(\cos(x^{2})\) is the single most common chain-rule error, which the Caution below isolates.
\end{envsolution}
\end{workedexample}
\begin{envcaution}{Caution}{}{}
The most common mistake with the chain rule is to differentiate the outer function and forget the inner derivative, writing \(\dfrac{d}{dx}\sin(g(x)) = \cos(g(x))\) instead of the correct \(\cos(g(x))\cdot g'(x)\). The inner factor \(g'(x)\) — the \(2x\) in \(\sin(x^{2})\) above — is never optional.
\end{envcaution}
\begin{workedexample}
\begin{envexample}{Example}{ex:3.5}{}
Differentiate \(\sqrt{1 + \sin^{2} x}\).
\end{envexample}
\begin{envsolution}{Solution}{}{}
This is a composition three layers deep: the square root of \(1 + \sin^{2} x\), in which \(\sin^{2} x = (\sin x)^{2}\) is itself a composition. Peeling from the outside in, the outer function is \(\sqrt{\,\cdot\,}\) applied to \(u = 1 + \sin^{2} x\); inside, \(\sin^{2} x\) is the square of \(v = \sin x\). Differentiating one layer at a time,

\[ \begin{aligned}
          \frac{d}{dx}\sqrt{1 + \sin^{2} x}
            &= \frac{1}{2\sqrt{1 + \sin^{2} x}}\cdot \frac{d}{dx}\bigl(1 + \sin^{2} x\bigr) \\[2pt]
            &= \frac{1}{2\sqrt{1 + \sin^{2} x}}\cdot 2\sin x\cdot \cos x \\[2pt]
            &= \frac{\sin x \cos x}{\sqrt{1 + \sin^{2} x}}.
        \end{aligned} \]

Each layer contributes exactly one factor: \(\tfrac12 u^{-1/2}\) from the square root, \(2\sin x\) from squaring, and \(\cos x\) from the sine. Their product is the derivative. The slopes multiply down the chain, just as Figure \ref{fig:3.5} pictures.
\end{envsolution}
\end{workedexample}
\begin{workedexample}
\begin{envexample}{Example}{ex:3.6}{}
Differentiate \(\sqrt{\dfrac{x - 1}{x + 2}}\) for \(x > 1\).
\end{envexample}
\begin{envsolution}{Solution}{}{}
The outermost operation is the square root, so by Strategy \ref{strat:3.1} take \(f(u) = \sqrt{u} = u^{1/2}\) with inner function \(u = \dfrac{x - 1}{x + 2}\). The chain rule gives

\[ \frac{d}{dx}\sqrt{\frac{x - 1}{x + 2}} = \frac{1}{2\sqrt{(x - 1)/(x + 2)}}\cdot\frac{d}{dx}\!\left(\frac{x - 1}{x + 2}\right). \]

The inner derivative is now a quotient-rule computation (§2.5):

\[ \frac{d}{dx}\!\left(\frac{x - 1}{x + 2}\right) = \frac{(x + 2)\cdot 1 - (x - 1)\cdot 1}{(x + 2)^{2}} = \frac{3}{(x + 2)^{2}}. \]

Substituting, and writing \(\dfrac{1}{2\sqrt{(x - 1)/(x + 2)}} = \dfrac{\sqrt{x + 2}}{2\sqrt{x - 1}}\),

\[ \frac{d}{dx}\sqrt{\frac{x - 1}{x + 2}} = \frac{\sqrt{x + 2}}{2\sqrt{x - 1}}\cdot\frac{3}{(x + 2)^{2}} = \frac{3}{2\sqrt{x - 1}\,(x + 2)^{3/2}} \qquad (x > 1). \]

The chain rule does not finish the computation by itself. It requires the derivative of the inside, and here that derivative is itself a quotient-rule computation.
\end{envsolution}
\end{workedexample}
\begin{workedexample}
\begin{envexample}{Example}{ex:3.7}{}
Differentiate \(y = (1 + x^{2})\cos^{2} x\).
\end{envexample}
\begin{envsolution}{Solution}{}{}
The outermost operation is a product, so begin with the product rule (§2.5):

\[ y' = (2x)\cos^{2} x + (1 + x^{2})\,\frac{d}{dx}\bigl(\cos^{2} x\bigr). \]

The remaining derivative needs the chain rule, since \(\cos^{2} x = [\cos x]^{2}\) is a composition with outer function \((\,\cdot\,)^{2}\) and inner function \(\cos x\):

\[ \frac{d}{dx}\bigl(\cos^{2} x\bigr) = 2\cos x\cdot(-\sin x) = -2\sin x\cos x. \]

Therefore

\[ y' = 2x\cos^{2} x - 2(1 + x^{2})\sin x\cos x. \]

The two rules cooperate: the product rule splits the work into two terms, and the chain rule supplies the slope factor inside the second. Neither alone differentiates this function.
\end{envsolution}
\end{workedexample}
\begin{workedexample}
\begin{envexample}{Example}{ex:3.8}{}
In a harbour, the amount of kelp \(K\) depends on the urchin population \(U\), since urchins eat kelp, and \(U\) in turn depends on the otter population \(O\), since otters eat urchins. Use the chain rule to decide whether \(\dfrac{dK}{dO}\) is positive or negative.
\end{envexample}
\begin{envsolution}{Solution}{}{}
Determine the sign of each link first. More urchins means less kelp, so \(K\) decreases as \(U\) increases: \(\dfrac{dK}{dU} < 0\). More otters means fewer urchins, so \(\dfrac{dU}{dO} < 0\). The chain rule in Leibniz form (Remark \ref{rem:3.2}) chains the two rates:

\[ \frac{dK}{dO} = \frac{dK}{dU}\cdot\frac{dU}{dO}. \]

A product of two negative numbers is positive, so \(\dfrac{dK}{dO} > 0\): more otters ultimately means more kelp. This is the formal version of the two-step argument that otters thin out the urchins and that fewer urchins let the kelp recover. Each \(dy/du\) factor carries the sign of one link, and the signs multiply along the chain.
\end{envsolution}
\end{workedexample}
With the chain rule in hand the differentiation toolkit is essentially complete: any function assembled by composing the elementary functions of Chapters 2 and 3 can now be differentiated mechanically. In the next section the rule does something new. It gives us derivatives we could not otherwise reach, those of the inverse functions \(\ln x\), \(\arcsin x\), and \(\arctan x\), and, through logarithmic differentiation, those of expressions such as \(x^{x}\).


\sechead{3.3}{Applications of the Chain Rule}
Sections 3.1 and 3.2 built the machinery; this section uses it. The chain rule (Theorem \ref{thm:3.3}) was advertised as the tool that completes the differentiation toolkit, and here we see why: it reaches a whole family of derivatives that no rule of Chapter 2 can differentiate directly. We have no formula for \(\ln x\) to differentiate. It is known only as the function that undoes \(e^{x}\). The expression \(x^{x}\) carries its variable in \emph{both} the base and the exponent, so neither the power rule nor the rule for \(e^{x}\) applies to it. And \(\arcsin\) has no elementary formula at all. Each of these functions is defined only indirectly.

Yet each one satisfies an \emph{identity} assembled from functions we already know how to differentiate: \(e^{\ln x} = x\), or \(\ln(x^{x}) = x\ln x\), or \(\arcsin(\sin x) = x\) for \(x \in \bigl[-\tfrac{\pi}{2}, \tfrac{\pi}{2}\bigr]\). That is the whole idea of this section, used over and over: \emph{when no rule differentiates a function directly, differentiate an identity it satisfies, then solve for the derivative you want.} It is the chain rule run backward. In Strategy \ref{strat:3.1} we start from known derivatives and apply the chain rule to a composition. Here we apply it to a known relation and solve for the one derivative we do not yet have. The same method handles the inverse functions \(\ln\), \(\arcsin\), \(\arccos\), and \(\arctan\), and, through the closely related trick of \emph{logarithmic differentiation}, awkward powers like \(x^{x}\).

\subsechead{The logarithm and logarithmic differentiation}
We begin with the natural logarithm. Throughout this section we use \(\ln x\) informally, as the inverse of the exponential function. That is, \(\ln x\) is the number whose exponential is \(x\), so that \(e^{\ln x} = x\) for every \(x > 0\). This relation is all we need to find its derivative. The rigorous construction of \(\ln x\) as the inverse of \(e^{x}\), together with a proof that it is differentiable, is deferred to Chapter 4 (§4.5), which re-derives the formula below from the ground up.

\begin{workedexample}
\begin{envexample}{Example}{ex:3.9}{}
Show that \(\dfrac{d}{dx}\ln x = \dfrac{1}{x}\) for \(x > 0\).
\end{envexample}
\begin{envsolution}{Solution}{}{}
Start from the identity that defines the logarithm, \(e^{\ln x} = x\), and differentiate both sides with respect to \(x\). The left side is a composition, with outer function \(\exp\) and inner function \(\ln x\). By the chain rule (Theorem \ref{thm:3.3}), using \(\dfrac{d}{du}e^{u} = e^{u}\) (Chapter 2, §2.4, Theorem 2.5),

\[ \frac{d}{dx}\,e^{\ln x} = e^{\ln x}\cdot\frac{d}{dx}\ln x = x\cdot\frac{d}{dx}\ln x, \]

where the last step uses \(e^{\ln x} = x\) once more. The right side is simply \(\dfrac{d}{dx}x = 1\). Equating the two,

\[ x\cdot\frac{d}{dx}\ln x = 1, \qquad\text{so}\qquad \frac{d}{dx}\ln x = \frac{1}{x} \quad (x > 0). \]

The unknown derivative sat inside a relation we could differentiate; one application of the chain rule and one division isolated it.
\end{envsolution}
\end{workedexample}
The derivative of the logarithm is more than one more entry in a table. It also makes a general technique possible. When a function carries its variable in the exponent, or is a long product and quotient of many factors, taking its logarithm first turns powers into products and products into sums, all of which we can already differentiate. Applying \(\dfrac{d}{dx}\ln\) to both sides and then solving is called \emph{logarithmic differentiation}.

\begin{envstrategy}{Strategy}{strat:3.2}{Logarithmic differentiation}
To differentiate a complicated positive expression \(y = f(x)\), especially one with the variable in an exponent:

\begin{steps}
  \item Take natural logarithms of both sides: \(\ln y = \ln f(x)\).
  \item Simplify the right side with the laws of logarithms — \(\ln(ab) = \ln a + \ln b\), \(\ln\frac{a}{b} = \ln a - \ln b\), \(\ln(a^{b}) = b\ln a\) — until only sums of simple terms remain.
  \item Differentiate both sides with respect to \(x\). Since \(y\) is a function of \(x\), the left side is \(\dfrac{d}{dx}\ln y = \dfrac{y'}{y}\) by the chain rule.
  \item Solve for \(y'\) and substitute \(y = f(x)\) back in: \(y' = f(x)\cdot\dfrac{d}{dx}\bigl[\ln f(x)\bigr]\).
\end{steps}
In particular this differentiates any power-of-a-function \(f(x)^{g(x)}\) (with \(f(x) > 0\)): its logarithm is \(g(x)\ln f(x)\); differentiate that, then multiply back by \(f(x)^{g(x)}\). The next example is the cleanest case.
\end{envstrategy}
\begin{workedexample}
\begin{envexample}{Example}{ex:3.10}{}
Differentiate \(W(x) = x^{x}\) for \(x > 0\).
\end{envexample}
\begin{envsolution}{Solution}{}{}
There is no rule for a variable raised to a variable power, so we take logarithms. With \(W = x^{x} > 0\),

\[ \ln W = \ln\bigl(x^{x}\bigr) = x\ln x. \]

Differentiate both sides with respect to \(x\). On the left, the chain rule gives \(\dfrac{d}{dx}\ln W = \dfrac{W'}{W}\). On the right, the product rule (Chapter 2, §2.5, Theorem 2.6) together with \(\dfrac{d}{dx}\ln x = \dfrac{1}{x}\) (Example \ref{ex:3.9}) gives

\[ \frac{d}{dx}\bigl(x\ln x\bigr) = 1\cdot\ln x + x\cdot\frac{1}{x} = \ln x + 1. \]

Therefore \(\dfrac{W'}{W} = \ln x + 1\), and solving for \(W'\),

\[ \frac{d}{dx}x^{x} = W'(x) = W\,(\ln x + 1) = (1 + \ln x)\,x^{x} \quad (x > 0). \]
\end{envsolution}
\end{workedexample}
\begin{workedexample}
\begin{envexample}{Example}{ex:3.11}{}
Differentiate \(2^{x}\) for \(x \in \mathbb{R}\).
\end{envexample}
\begin{envsolution}{Solution}{}{}
Rewrite the base through \(e\). Since \(2 = e^{\ln 2}\), we have \(2^{x} = e^{x\ln 2}\), a composition with inner function \(x\ln 2\) (a constant times \(x\)) and outer function \(\exp\). By the chain rule, with \(\dfrac{d}{du}e^{u} = e^{u}\),

\[ \frac{d}{dx}2^{x} = \frac{d}{dx}e^{x\ln 2} = e^{x\ln 2}\cdot\ln 2 = 2^{x}\ln 2. \]

The same rewrite gives \(\dfrac{d}{dx}a^{x} = a^{x}\ln a\) for any base \(a > 0\). This is also why \(e\) is the \emph{natural} base: because \(\ln e = 1\), the extra factor disappears, and \(e^{x}\) alone among exponentials is its own derivative.
\end{envsolution}
\end{workedexample}
\begin{workedexample}
\begin{envexample}{Example}{ex:3.12}{}
Differentiate \(x\ln x - x\) for \(x > 0\).
\end{envexample}
\begin{envsolution}{Solution}{}{}
We met the first term inside Example \ref{ex:3.10}: by the product rule and \(\dfrac{d}{dx}\ln x = \dfrac{1}{x}\), we have \(\dfrac{d}{dx}(x\ln x) = \ln x + 1\). The second term has derivative \(1\). Hence

\[ \frac{d}{dx}\bigl(x\ln x - x\bigr) = (\ln x + 1) - 1 = \ln x \quad (x > 0). \]

The derivative is exactly \(\ln x\), so \(x\ln x - x\) is an \emph{antiderivative} of \(\ln x\), a fact worth remembering for when we turn to integration.
\end{envsolution}
\end{workedexample}
\begin{workedexample}
\begin{envexample}{Example}{ex:3.13}{}
Differentiate \(y = (x + 1)(x^{2} + 1)^{2}(x^{3} + 1)^{3}\) for \(x > 0\) by logarithmic differentiation.
\end{envexample}
\begin{envsolution}{Solution}{}{}
Differentiating this product directly would mean nesting the product rule again and again. Instead take logarithms. Since \(y > 0\),

\[ \ln y = \ln\!\bigl[(x + 1)(x^{2} + 1)^{2}(x^{3} + 1)^{3}\bigr]. \]

The laws of logarithms (Strategy \ref{strat:3.2}, step 2) turn the product into a sum and bring each exponent out in front:

\[ \ln y = \ln(x + 1) + 2\ln(x^{2} + 1) + 3\ln(x^{3} + 1). \]

Differentiate both sides with respect to \(x\). On the left, \(\dfrac{d}{dx}\ln y = \dfrac{y'}{y}\) by the chain rule; on the right, \(\dfrac{d}{dx}\ln u = \dfrac{u'}{u}\) gives one fraction per factor:

\[ \frac{y'}{y} = \frac{1}{x + 1} + \frac{2(2x)}{x^{2} + 1} + \frac{3(3x^{2})}{x^{3} + 1} = \frac{1}{x + 1} + \frac{4x}{x^{2} + 1} + \frac{9x^{2}}{x^{3} + 1}. \]

Solving for \(y'\) and substituting \(y\) back,

\[ y' = (x + 1)(x^{2} + 1)^{2}(x^{3} + 1)^{3}\left(\frac{1}{x + 1} + \frac{4x}{x^{2} + 1} + \frac{9x^{2}}{x^{3} + 1}\right). \]

Each factor of the original product contributes exactly one term to the bracket: logarithmic differentiation replaces one hard product-rule computation with several easy logarithmic derivatives added together.
\end{envsolution}
\end{workedexample}
\subsechead{Derivatives of the inverse trigonometric functions}
We found the derivative of the logarithm by differentiating the relation \(e^{\ln x} = x\) that expresses \(\ln\) as the inverse of \(e^{x}\). The same method works for the inverse trigonometric functions, only now the identity comes from trigonometry. Each inverse function composes with the function it inverts to give an identity, and differentiating that identity by the chain rule isolates the derivative we are after.

\begin{workedexample}
\begin{envexample}{Example}{ex:3.14}{}
Show that \(\dfrac{d}{dy}\arcsin y = \dfrac{1}{\sqrt{1 - y^{2}}}\) for \(y \in (-1, 1)\).
\end{envexample}
\begin{envsolution}{Solution}{}{}
Recall from Chapter 1 that \(\arcsin\) is the inverse of sine restricted to the branch \(\bigl[-\tfrac{\pi}{2}, \tfrac{\pi}{2}\bigr]\): for \(-1 \le y \le 1\), \(\arcsin y\) is the unique angle \(x\) in that interval with \(\sin x = y\). On this branch the two functions undo each other, so in particular \(\arcsin(\sin x) = x\).

Differentiate this identity with respect to \(x\). Reading \(x = \arcsin(\sin x)\) as a composition and applying the chain rule (Theorem \ref{thm:3.3}), with \(\dfrac{d}{dx}\sin x = \cos x\) (Theorem \ref{thm:3.1}),

\[ 1 = \frac{d}{dx}\arcsin(\sin x) = (\arcsin)'(\sin x)\cdot\cos x, \]

so \((\arcsin)'(\sin x) = \dfrac{1}{\cos x}\), valid where \(\cos x \neq 0\), that is, on the \emph{open} interior \(x \in \bigl(-\tfrac{\pi}{2}, \tfrac{\pi}{2}\bigr)\), where \(\cos x > 0\). Writing \(y = \sin x\), this interior corresponds exactly to \(y \in (-1, 1)\), and there \(\cos x = \sqrt{1 - \sin^{2} x} = \sqrt{1 - y^{2}}\), the positive root selected by \(\cos x > 0\). Therefore

\[ \frac{d}{dy}\arcsin y = \frac{1}{\sqrt{1 - y^{2}}} \quad (y \in (-1, 1)). \]
\end{envsolution}
\end{workedexample}
\begin{envcaution}{Caution}{}{}
The formula \(\dfrac{1}{\sqrt{1 - y^{2}}}\) is valid only on the \emph{open} interval \((-1, 1)\). At the endpoints \(y = \pm 1\) the denominator vanishes: the graph of \(\arcsin\) (and of \(\arccos\) below) has a vertical tangent there (Figure \ref{fig:3.7}), so the derivative does not exist. Report these derivatives on the open interval, never the closed one.
\end{envcaution}
\begin{figureblock}
  \includegraphics[width=64.03mm]{ch03/figs/arcsin-vertical-tangent}
\figcaption{fig:3.7}{Endpoint behavior of \(\arcsin y\): the tangent becomes vertical at \(y = \pm 1\), so the derivative formula applies only on \((-1, 1)\).}
\end{figureblock}
\begin{workedexample}
\begin{envexample}{Example}{ex:3.15}{}
Show that \(\dfrac{d}{dy}\arccos y = \dfrac{-1}{\sqrt{1 - y^{2}}}\) for \(y \in (-1, 1)\).
\end{envexample}
\begin{envsolution}{Solution}{}{}
The method is exactly that of Example \ref{ex:3.14}, now applied to cosine. Recall that \(\arccos\) is the inverse of cosine on the branch \([0, \pi]\), so \(\arccos(\cos x) = x\) there. Differentiating with respect to \(x\), with \(\dfrac{d}{dx}\cos x = -\sin x\) (Theorem \ref{thm:3.2}),

\[ 1 = (\arccos)'(\cos x)\cdot(-\sin x). \]

Writing \(y = \cos x\): on the \emph{open} interior \(x \in (0, \pi)\) we have \(\sin x > 0\), so the division by \(-\sin x\) is legitimate and \(\sin x = \sqrt{1 - \cos^{2} x} = \sqrt{1 - y^{2}}\); this interior corresponds to \(y \in (-1, 1)\). Hence

\[ \frac{d}{dy}\arccos y = \frac{1}{-\sin x} = \frac{-1}{\sqrt{1 - y^{2}}} \quad (y \in (-1, 1)). \]

The only change from \(\arcsin\) is the minus sign, inherited from \(\cos' = -\sin\); the open-interval restriction of the Caution applies here unchanged.
\end{envsolution}
\end{workedexample}
\begin{workedexample}
\begin{envexample}{Example}{ex:3.16}{}
Show that \(\dfrac{d}{dy}\arctan y = \dfrac{1}{1 + y^{2}}\) for all \(y \in \mathbb{R}\).
\end{envexample}
\begin{envsolution}{Solution}{}{}
Once more the same method. Recall that \(\arctan\) is the inverse of tangent on the branch \(\bigl(-\tfrac{\pi}{2}, \tfrac{\pi}{2}\bigr)\), so \(\arctan(\tan x) = x\). Differentiating with respect to \(x\), and using \(\dfrac{d}{dx}\tan x = \sec^{2} x\) (Example \ref{ex:3.2}),

\[ 1 = (\arctan)'(\tan x)\cdot\sec^{2} x, \qquad\text{so}\qquad (\arctan)'(\tan x) = \frac{1}{\sec^{2} x}. \]

Now write \(y = \tan x\) and use the Pythagorean identity \(\sec^{2} x = 1 + \tan^{2} x = 1 + y^{2}\):

\[ \frac{d}{dy}\arctan y = \frac{1}{1 + y^{2}}. \]

Here there is no domain restriction to worry about: as \(x\) ranges over \(\bigl(-\tfrac{\pi}{2}, \tfrac{\pi}{2}\bigr)\), \(y = \tan x\) ranges over all of \(\mathbb{R}\), and \(1 + y^{2}\) never vanishes. Unlike \(\arcsin\) and \(\arccos\), the arctangent is differentiable everywhere.
\end{envsolution}
\end{workedexample}
\begin{envremark}{Remark}{rem:3.3}{One lever for many derivatives}
Look back at what just happened. We never differentiated \(\ln x\), \(x^{x}\), or an inverse trigonometric function \emph{directly}. There was no rule to do so. In every case we found a known identity the function satisfies, differentiated it by the chain rule, and solved for the one unknown. The identities were \(e^{\ln x} = x\), \(\ln(x^{x}) = x\ln x\), \(\arcsin(\sin x) = x\), \(\arccos(\cos x) = x\), and \(\arctan(\tan x) = x\), each on the branch fixed in its own example. The chain rule differentiates compositions, and applied to a known identity it also gives derivatives for which no direct rule exists.
\end{envremark}
\subsechead{Chapter summary}
This section adds the following derivatives to our running table:

\[ \frac{d}{dx}\ln x = \frac{1}{x}, \qquad \frac{d}{dx}x^{x} = (1 + \ln x)\,x^{x}, \qquad \frac{d}{dx}a^{x} = a^{x}\ln a, \]

\[ \begin{aligned}
    \frac{d}{dy}\arcsin y &= \frac{1}{\sqrt{1 - y^{2}}}, \\[2pt]
    \frac{d}{dy}\arccos y &= \frac{-1}{\sqrt{1 - y^{2}}}, \\[2pt]
    \frac{d}{dy}\arctan y &= \frac{1}{1 + y^{2}}.
  \end{aligned} \]

(The logarithmic and power formulas hold for \(x > 0\), the \(a^{x}\) formula for any base \(a > 0\) and all \(x\), and the \(\arcsin\) and \(\arccos\) formulas only on the open interval \(-1 < y < 1\); \(\arctan\) is differentiable for every \(y\).)

With these, Chapter 3 is complete. It opened by evaluating the one limit algebra could not reach, \(\lim_{\theta \to 0}\dfrac{\sin\theta}{\theta}\), with a piece of unit-circle geometry, and deriving \(\sin' = \cos\) and \(\cos' = -\sin\) (§3.1). It then proved the chain rule, the rule for compositions that Chapter 2 had been missing (§3.2). Here that rule was used to obtain the derivatives of the logarithmic, exponential, and inverse trigonometric functions (§3.3). Together with the algebraic rules of Chapter 2, we can now differentiate essentially every function built by combining and composing the elementary ones.

Two debts remain, and both were taken on informally. We have used \(\dfrac{d}{dx}e^{x} = e^{x}\) from the casual series argument of Chapter 2, and we have treated \(\ln x\) as the inverse of \(e^{x}\) without ever constructing it. Chapter 4 settles both: it builds \(e^{x}\) rigorously and re-derives its derivative (§4.3), and it proves the Mean Value Theorem (§4.4), the result that finally lets us define \(\ln x\) as the inverse of the strictly increasing exponential and recover \(\dfrac{d}{dx}\ln x = \dfrac{1}{x}\) rigorously (§4.5).
\end{document}
